\section{Conclusions at the Boundaries of the Frameworks}
\label{sec:conclusions}

\noindent
The paper terminates as the prior papers of the series have, not in a single conclusion that would gather the frameworks into a possessed whole, but in plural conclusions, each stated from within a framework to the boundary of what that framework establishes. The plurality is required by the method: a flourishing that no single framework possesses cannot be concluded in the terms of any one, and the dialogue among the frameworks, each speaking to its boundary, is the only form a conclusion may take without installing one framework as the metalanguage the series denies.

\subsection{From positivism}
\label{sec:concl-positivism}

Positivism establishes that flourishing has measurable indicators, that these can be brought to validated and comparable instruments, and that the structured measures so obtained are a public and accumulable gain over private appraisal. It establishes the recyclability, the skew of return, and the rates of accomplished and foreclosed undertaking as structured proxies that make visible what aggregate satisfaction conceals. Positivism does not establish that its measures reach the value they proxy, which is unsymbolised and beyond their structuring; it cannot distinguish, from a satisfaction score alone, a flourishing spiral from a forged one; and it reaches its boundary at the point where the structured proxy must be read against the value it does not contain, which it consigns to reflection and to the basal appraisal it cannot replace.

\subsection{From evidentialism}
\label{sec:concl-evidentialism}

Evidentialism establishes that belief about flourishing, and the practice founded on it, ought to be proportioned to evidence and revised as evidence revises, and that this proportioning operates not only in the explicit weighing of studies but in the continuous belief updating that underlies all appraisal. It establishes the supported practices a cultivation may draw upon, and the form of the hypothesis test by which the justice of a value cycle might be assessed. Evidentialism does not establish what the singular shared life requires beyond the average its evidence reports, and it is mute where the value at issue cannot be brought to evidence at all; it reaches its boundary at the unsymbolised appraisal that is its own deep form and yet exceeds any evidence it can marshal.

\subsection{From normative theory}
\label{sec:concl-normative}

Normative theory establishes that flourishing is the substantive freedom of each party to unfold what that party has reason to value, that this unfolding is plural and not aggregate, and that a just flourishing requires the return of value to its generators and the foreclosure of no party's unfolding. It establishes the normative specification against which the paper's evaluation is conducted. Normative theory does not establish that its specification is realised in a given life, which is a matter of practice and assessment under actual conditions; and it does not reach the lived and unsymbolised meaning in which the specification has its concrete sense for a particular couple, which it consigns to the phenomenological.

\subsection{From virtue ethics}
\label{sec:concl-virtue}

Virtue ethics establishes that flourishing is activity in accordance with cultivated excellence, that such excellence is formed by habituation in the concrete fabric of a shared life, and that the cultivation of disposition, and not the legislation of it, is the just form of this formation. It establishes the connection between the cognition of flourishing and its practice that the other frameworks leave implicit. Virtue ethics does not establish a procedure for the formation of the right dispositions, since a procedure applied from outside the joint practice would be the legislation it forbids; it reaches its boundary at the point where the cultivation of a shared disposition must proceed without a rule, in the unguaranteeable practice the prior paper described.

\subsection{From phenomenology}
\label{sec:concl-phenomenology}

Phenomenology establishes that flourishing is lived and interpreted before it is measured, that its givenness in the felt quality and the apprehended meaning of a shared life is the matter to which every measure is answerable, and that the basal appraisal in which a relation is continuously evaluated is conducted in this lived and unsymbolised register. It establishes the limit against which the structured measures are to be read. Phenomenology does not establish anything public, comparable, or accumulable, since what it describes resists exactly the structuring that would make it so; it reaches its boundary at the point where the lived meaning of a shared life must be communicated, assessed, or accumulated, which it consigns to the structured proxies it cannot itself supply.

\subsection{From the cognitive and dynamical account}
\label{sec:concl-dynamical}

The cognitive and dynamical account establishes that appraisal is, at its basis, a continuous inference upon a partially observable state, conducted in unsymbolised value; that the relation is maintained by the persistence of this appraisal and the practice taken upon it, so that the cessation of the cycle is the relation's extinction; and that the cycle's regime is fixed by its persistence and the sign of its phase, with justice as the internal condition of a positive phase for all parties. It establishes the lifecycle as a real cycle and the three regimes as its possible fates. The account does not establish a fixed model of the cycle, since the value function of the relational decision process is itself revised by the cycle, in the manner the prior paper analysed as the absence of a fixed meta-value-grammar; it reaches its boundary at the reflexive reconfiguration that no fixed model captures, and consigns the matter, as the prior paper did, to the open question of a formalism whose own ground is not fixed.

\subsection{The absence of a single conclusion}
\label{sec:concl-none}

These conclusions do not sum. Each framework, developed to its boundary, establishes what it establishes and consigns to another what it cannot reach: positivism consigns the unsymbolised remainder to phenomenology and reflection, normative theory consigns the lived meaning to phenomenology and the realisation to practice, virtue ethics consigns the formation without a rule to the unguaranteeable practice, phenomenology consigns the public and comparable to the structured proxy, and the dynamical account consigns the reflexive remainder to the open formal question. The flourishing of a shared life is the object none of them possesses and all of them partially determine, and the operational concession governs the whole: where the paper assessed, it assessed a structured proxy for a value it did not reach, and it bore the incompleteness rather than concealing it.

\begin{claim}[The absence of a single conclusion as the adequate conclusion]
There is no single conclusion that gathers the frameworks into a possessed account of flourishing, not for want of one but because flourishing is not possessed by any framework and is in production across them. The plural conclusions are partial and mutually irreducible of necessity, and their dialogue, each framework speaking to its boundary and consigning what it cannot reach, is itself the polyphonic and unconcluded form that a just eudaimonia, cultivated and not constructed, assessed in proxy and not possessed, sustained as an open spiral and not secured as a state, requires of any account faithful to it. That the paper does not conclude in a single account is its fidelity to an object that is cultivated, processual, and unguaranteeable: the account, like the flourishing it treats, is sustained as an open movement and not secured as a possessed result.
\end{claim}

\noindent
The flourishing of a shared life is, in the end, neither measured nor possessed but lived, cultivated, and continually reappraised, in a cycle whose persistence is its life, whose phase is its rising or its decline, and whose justice is the condition of its rising for both. The frameworks illuminate it severally and from their boundaries; the cycle turns; and the good of a shared life is generated, where it is generated, in the turning, and not in any account of it, this one included.
